The Brazilian National High School Examination (ENEM) is the primary pathway to
higher education in Brazil. However, recent results indicate a persistent
decline in students' mathematics performance. International assessments, such as
the 2022 Programme for International Student Assessment (PISA), corroborate this
concerning scenario by showing that a substantial proportion of Brazilian
students face difficulties in mathematics, while only a small percentage achieve
high levels of proficiency. Previous studies indicate that these difficulties
arise not only from gaps in mathematical knowledge but also from inadequate
cognitive strategies, limited comprehension of problem statements, and low
levels of student engagement. Although the Brazilian National Common Curricular
Base (BNCC) recognizes Computational Thinking (CT) as a fundamental competency
for problem-solving, few educational tools effectively integrate CT principles
with engaging instructional strategies. This paper presents Cognitivo, a
gamified educational software designed around the four pillars of CT, namely decomposition, abstraction, pattern recognition, and algorithm
design, to support students preparing for the mathematics section of the ENEM.
By combining gamification with CT principles, the proposed
approach aims to promote student engagement, logical reasoning, and autonomous
problem-solving in contextualized mathematical tasks. Consequently, Cognitivo
seeks to address key challenges in mathematics education while providing an
educational tool aligned with the competencies established by the BNCC.


\keywords{Computational Thinking; Gamification; Engagement; Problem Solving; ENEM; Mathematics Education; }